% =====================================================================
%  Cause-for-Concern (.cfc)
%  A Domain-Specific Language for Disease Diagnosis and Therapy by
%  Coherence Localization and Gap-Closure.
%
%  Self-contained. Every object is defined from first principles; the
%  reader is assumed to know undergraduate real analysis, elementary
%  graph theory, and the standard vocabulary of programming-language
%  semantics. No companion manuscript is required.
% =====================================================================
\documentclass[11pt,a4paper]{article}

\usepackage[utf8]{inputenc}
\usepackage[T1]{fontenc}
\usepackage{lmodern}
\usepackage[a4paper,margin=2.4cm]{geometry}
\usepackage{amsmath,amssymb,amsthm,mathtools}
\usepackage{bm}
\usepackage{mathrsfs}
\usepackage{booktabs}
\usepackage{array}
\usepackage{enumitem}
\usepackage{microtype}
\usepackage{xcolor}
\usepackage{listings}
\usepackage{algorithm}
\usepackage{algpseudocode}
\usepackage[round,sort&compress]{natbib}
\usepackage[colorlinks=true,linkcolor=blue!45!black,
            citecolor=blue!45!black,urlcolor=blue!45!black]{hyperref}
\usepackage{cleveref}

% ---- Theorem environments ----
\newtheoremstyle{cfcthm}{\topsep}{\topsep}{\itshape}{}{\bfseries}{.}{.5em}{}
\theoremstyle{cfcthm}
\newtheorem{theorem}{Theorem}[section]
\newtheorem{lemma}[theorem]{Lemma}
\newtheorem{proposition}[theorem]{Proposition}
\newtheorem{corollary}[theorem]{Corollary}
\theoremstyle{definition}
\newtheorem{definition}[theorem]{Definition}
\newtheorem{assumption}{Assumption}
\newtheorem{principle}[theorem]{Principle}
\newtheorem{rulename}[theorem]{Rule}
\newtheorem{problem}[theorem]{Problem}
\newtheorem{example}[theorem]{Example}
\theoremstyle{remark}
\newtheorem{remark}[theorem]{Remark}
\newtheorem{notation}[theorem]{Notation}

\crefname{assumption}{Assumption}{Assumptions}
\Crefname{assumption}{Assumption}{Assumptions}
\crefname{rulename}{Rule}{Rules}
\Crefname{rulename}{Rule}{Rules}
\crefname{principle}{Principle}{Principles}
\Crefname{principle}{Principle}{Principles}

% ---- Shorthand ----
\newcommand{\R}{\mathbb{R}}
\newcommand{\N}{\mathbb{N}}
\newcommand{\Rp}{\R_{\ge 0}}
\newcommand{\Man}{\mathcal{M}}          % state manifold
\newcommand{\st}{x}                     % state
\newcommand{\coh}{x^{\star}}            % coherent state
\newcommand{\Coh}{\Man^{\star}}         % coherent set
\newcommand{\flo}{\beta}                % floor
\newcommand{\cell}{\mathcal{C}}         % cell
\newcommand{\probe}{\pi}                % probe
\newcommand{\Probe}{\Pi}                % probe space
\newcommand{\trace}{\tau}               % trace
\newcommand{\Rd}{\rho}                  % readout map
\newcommand{\fwd}{\Phi}                 % forward model / trace map
\newcommand{\loc}{\Lambda}              % localization
\newcommand{\gap}{\mathsf{g}}           % gap
\newcommand{\basis}{\mathcal{B}}        % intervention basis
\newcommand{\intv}{u}                   % intervention
\newcommand{\dose}{c}                   % dose vector
\newcommand{\cfc}{\mathcal{L}}          % the language
\newcommand{\ecoh}{\eta}                % coherence index
\newcommand{\dvec}{\mathbf{D}}          % disease vector
\newcommand{\Dom}{\Omega}               % oscillator index set
\newcommand{\rec}{M}                    % committed-measurement count
\newcommand{\abs}[1]{\lvert#1\rvert}
\newcommand{\norm}[1]{\lVert#1\rVert}
\newcommand{\set}[1]{\{#1\}}
\newcommand{\eps}{\varepsilon}
\newcommand{\diam}{\operatorname{diam}}
\newcommand{\dist}{\operatorname{dist}}
\newcommand{\supp}{\operatorname{supp}}
\DeclareMathOperator*{\argmin}{arg\,min}
\DeclareMathOperator*{\argmax}{arg\,max}
\newcommand{\stepto}{\longrightarrow}
\newcommand{\bnf}{\;::=\;}
\newcommand{\alt}{\;\mid\;}
\newcommand{\cf}[1]{\texttt{#1}}
\newcommand{\nt}[1]{\ensuremath{\langle\textsf{#1}\rangle}}

% ---- .cfc listing style ----
\definecolor{cfckw}{RGB}{30,64,175}
\definecolor{cfccom}{RGB}{22,101,52}
\definecolor{cfcnum}{RGB}{180,83,9}
\definecolor{cfcstr}{RGB}{159,18,57}
\definecolor{cfcbg}{RGB}{250,250,250}
\lstdefinelanguage{cfc}{
  morekeywords={floor,patient,probe,panel,measure,localize,cell,gap,
    close,toward,coherent,using,basis,until,resolved,contested,decline,
    concern,when,emit,assert,yield,let,in,as,with,import,module,export,
    monitor,every,observe,against,reference},
  morekeywords=[2]{Coherence,DiseaseVector,Cell,Gap,Probe,Panel,
    Intervention,Dose,Reading,Verdict,Concern},
  sensitive=true,
  morecomment=[l]{--},
  morestring=[b]",
  keywordstyle=\color{cfckw}\bfseries,
  keywordstyle=[2]\color{cfcnum},
  commentstyle=\color{cfccom}\itshape,
  stringstyle=\color{cfcstr},
  basicstyle=\ttfamily\small,
  backgroundcolor=\color{cfcbg},
  frame=single,framesep=4pt,xleftmargin=6pt,xrightmargin=4pt,
  breaklines=true,showstringspaces=false,columns=fullflexible,tabsize=2
}
\lstset{language=cfc}

\title{\textbf{Cause-for-Concern: A Domain-Specific Language for\\[2pt]
Disease Diagnosis and Therapy by Coherence Localization\\
and Gap-Closure}}

\author{%
Kundai Farai Sachikonye\\
\small Technical University of Munich\\
\small Chair of Experimental Bioinformatics\\
\small AIMe Registry for Artificial Intelligence\\
\small \texttt{kundai.sachikonye@wzw.tum.de}}

\date{\today}

\begin{document}
\maketitle

% =====================================================================
\begin{abstract}
\noindent
We specify \emph{Cause-for-Concern} (\cf{cfc}; source extension \cf{.cfc}),
a domain-specific language for clinical reasoning whose semantics dispenses
with the two classical primitives of medicine---the patient's narrative
report and the named disease drawn from a catalogue---and replaces them with
a single primitive: the \emph{characterised perturbation} of a patient's
physiological state, whose fate is computed against a forward model and read
back as a measurement. The language is self-contained and rests on five facts
we prove rather than assume. \textbf{(i)}~A finite observer of a physiological
system cannot resolve its state to a point; a strictly positive resolution
floor $\flo>0$ forces every determined state to be a \emph{cell} of diameter
$\ge\flo$, so a diagnosis is the identification of a cell, never a label
(\emph{Floor Theorem}). \textbf{(ii)}~Health is not a distinguished catalogue
entry but membership of a \emph{coherent set}, the states in which the
system's oscillator ensemble is internally phase-consistent, quantified by a
scalar coherence $\ecoh\in[0,1]$ and a disease vector
$\dvec\in[0,1]^{d}$; diagnosis is the \emph{localization} of the patient to a
cell, provably sufficient exactly when the probe-to-reading map is injective
down to the floor (\emph{Localization Theorem}). \textbf{(iii)}~Because each
physiological response is itself a substrate for the next, a probe can be made
to \emph{cascade}, so its terminal readout \emph{certifies} the trajectory
that produced it: a false account of the state cannot forge it, as it would
have to produce products that never formed (\emph{Certificate Theorem}), and
the correct halting rule is a \emph{rigidity} condition, not a magnitude
threshold. \textbf{(iv)}~Therapy is closure of the measured \emph{gap} to the
nearest coherent state by the sparsest combination of characterised
interventions selected \emph{by effect} and not by indication, solved as an
$\ell_{1}$ program and executed as a contraction whose fixed point is
coherence, with no disease ever named (\emph{Gap-Closure} and
\emph{Convergence Theorems}). \textbf{(v)}~Diagnosis and therapy are one
operation, since every characterised perturbation simultaneously informs
(its fate is a measurement) and displaces (its action is an intervention);
the loop halts precisely when the state is both coherent and rigidly known
(\emph{Wellness Theorem}). We give \cf{cfc}'s complete lexical grammar,
context-free grammar in EBNF, an \emph{accountability} type system in which no
value of zero residue is well-typed and no \cf{close} converges without a
localized cell, a small-step operational semantics in which evaluation is
measurement---mutating a monotone committed-measurement clock rather than
simulating---and proofs of type soundness (progress and preservation),
determinism, and scheduler soundness (a closed cell is never re-probed and a
convergent therapy is never starved). We specify complexity, a constructive
validation protocol requiring no clinical data, and a reference execution
model. The result is a deterministic, disease-free, report-free calculus in
which a patient's state is read from a self-certifying physiological program,
the therapeutic target is a measured gap, and clinical judgment is relocated
from precedent-retrieval to the specification of \emph{intent} over a derived,
explorable state.
\end{abstract}

\noindent\textbf{Keywords:} domain-specific language; disease diagnosis;
therapy; resolution floor; coherence; state localization; trajectory
certificate; gap-closure; sparse intervention; contraction to a fixed point;
accountability type system; operational semantics; measurement-as-evaluation.

\tableofcontents

% =====================================================================
\section{Introduction}
\label{sec:intro}
% =====================================================================

\subsection{Two primitives medicine cannot justify}

Clinical practice is organised around two primitives that its own
methodology cannot ground. The first is the \emph{patient's report}: a
narrative of symptoms, taken as the entry point to diagnosis. The second is
the \emph{named disease}: a label drawn from a finite nosological catalogue,
against which the report is matched. Both are load-bearing, and both are
fragile. A report is a lossy, subjective, and sometimes absent channel; a
catalogue is a historically contingent partition of a continuous
pathophysiology, with boundaries that shift between editions and that fit
poorly the many conditions which are combinations, prodromes, or genuinely
novel \citep{feinstein1967clinical,scheuermann2009toward,loscalzo2007network}.
When a clinician reasons ``the report resembles catalogue entry $k$, whose
standard therapy is $T_k$,'' the inference is retrieval of precedent, and its
validity is only as good as the catalogue's fit to the individual in front of
them.

This paper specifies a language whose execution never uses either primitive.
It replaces both with a single computational act---the administration of a
\emph{characterised perturbation} to the patient and the reading of its
fate---and shows that diagnosis, therapy, and the decision of when to stop all
reduce to operations on that one act. The language is called
\emph{Cause-for-Concern}, written \cf{cfc}, with source files carrying the
extension \cf{.cfc}. The name is literal: a \cf{cfc} program does not ask
``which disease?'' but ``is there a cause for concern, where, and by how
much?''---and answers with a region of state space and a measured
displacement, not a name.

\subsection{The one primitive: a characterised perturbation}

Let a patient's physiology be a point $\st$ in a state space $\Man$ (a full
description of the relevant oscillator ensemble: concentrations, activities,
conductances, phases). An ingested or administered compound of \emph{known}
composition is a controlled input to this system; its fate---the time-course
of products and physiological responses it induces---is a function of the
patient's state. If the forward model predicting that fate is valid across the
state space and not merely at a healthy reference, then the measured fate
\emph{locates} the patient in $\Man$. This single object, the characterised
perturbation, is read two ways: as a \emph{measurement} (its fate carries
information about the state) and as an \emph{intervention} (its action
displaces the state). \cf{cfc} exposes exactly this object, under the keyword
\cf{probe}, and derives everything else from it.

\subsection{Five facts, proved not assumed}

The language is not an interface pasted over heuristics; its every guarantee
is a theorem. Five facts organise the development, each established in the
body from stated premises.

\begin{enumerate}[leftmargin=1.5em,itemsep=3pt]
\item \textbf{The resolution floor} (\Cref{sec:floor}). A finite observer of a
physiological system cannot resolve its state to a point: there is a strictly
positive floor $\flo>0$ below which distinctions cannot be made. ``The state''
is therefore a \emph{cell} of diameter $\ge\flo$, and the goal of measurement
is not to name the state but to shrink the admissible region to a single cell.
\item \textbf{Diagnosis as localization} (\Cref{sec:coherence,sec:diagnosis}).
Health is membership of a \emph{coherent set} $\Coh$---the states in which the
oscillator ensemble is internally phase-consistent---quantified by a scalar
coherence $\ecoh\in[0,1]$ and a disease vector $\dvec$. A probe is
\emph{diagnostically sufficient} for a region iff its state-to-reading map is
injective down to the floor, in which case reading the probe pins the patient
to one cell. No catalogue enters.
\item \textbf{The trajectory certificate} (\Cref{sec:cascade}). Because a
physiological response is itself a substrate, a probe can \emph{cascade}: a
later product exists only if the earlier trajectory occurred as predicted, so
the readout is a \emph{certificate} of a path, not a snapshot. A false account
cannot forge it. The correct stopping rule is \emph{rigidity}: halt when every
account other than the true one is refuted by some observed product it cannot
explain.
\item \textbf{Therapy as gap-closure} (\Cref{sec:therapy}). The therapeutic
target is the vector from the localized state to the nearest coherent state.
Treatment is the sparsest combination of characterised interventions---drawn
from a basis indexed by \emph{effect}, not by indication---that drives that
vector to the floor, solved as an $\ell_1$ program and executed as a
contraction whose fixed point is coherence. The disease is never named.
\item \textbf{Diagnosis and therapy are one} (\Cref{sec:unity}). Every
characterised perturbation both measures and moves; the measure--act--remeasure
loop needs no external monitor and halts precisely when the state is both
coherent (gap below floor) and rigidly known (no rival account survives). This
double condition is the operative definition of \emph{wellness}.
\end{enumerate}

\subsection{What the language contributes}

On this foundation we build a full language, not a sketch. \Cref{sec:model}
fixes the semantic model. \Cref{sec:floor,sec:coherence,sec:diagnosis,%
sec:cascade,sec:therapy,sec:unity} prove the five facts and their supporting
lemmas. \Cref{sec:runtime} states the one substantive demand on the runtime
(that the forward model be \emph{derived}, not fitted) as a theorem about
extrapolation, because the safety of \emph{in silico} experimentation reduces
to it. \Cref{sec:lexical,sec:grammar} give the lexical and context-free
grammar; \Cref{sec:types} the accountability type system with proofs;
\Cref{sec:opsem} the small-step operational semantics with proofs of
monotonicity, determinism, and scheduler soundness; \Cref{sec:stdlib} the
standard library binding each verb to a semantic operation. \Cref{sec:examples}
gives worked programs, \Cref{sec:complexity} the complexity of each judgment,
\Cref{sec:validation} a constructive validation protocol requiring no clinical
data, and \Cref{sec:impl} the reference execution model. \Cref{sec:related}
situates the design; \Cref{sec:conclusion} concludes.

\subsection{Method and standard of rigour}

Every object below is a finite-dimensional dynamical system, a finite weighted
structure, or a construction on one, and every claim is stated with its
hypotheses and proved from them. Where a clinical reading is offered---``cell''
as differential diagnosis, ``gap'' as therapeutic deficit, ``decline'' as an
honest \emph{non liquet}---it is confined to a remark on which no theorem
depends. The grammar, type system, and semantics are the language as a
mathematical object; the present text, not any implementation, is normative.

% =====================================================================
\section{The Semantic Model}
\label{sec:model}
% =====================================================================

\cf{cfc} is defined over a single mathematical object---the physiological
state system---and three derived constructions: the coherence functional, the
cell, and the gap. We fix these before any syntax.

\subsection{The state system}

\begin{definition}[Physiological state system]
\label{def:system}
A \emph{physiological state system} is a tuple
$\Sigma=(\Man,F,\Probe,\Rd,\Dom)$ where:
\begin{itemize}[leftmargin=1.4em,itemsep=2pt]
\item $\Man\subseteq\Rp^{\,n}$ is the \emph{state space}: a point $\st\in\Man$
is a full description of the relevant physiology (metabolite concentrations,
enzyme activities, channel conductances, oscillator phases) at an instant. The
dimension $n$ is finite but may be large.
\item $F\colon\Man\times\Probe\to\R^{n}$ is the \emph{dynamics}:
$\dot\st=F(\st,\intv)$, with $\intv$ an input. $F$ is a kinetic law (mass
action, Michaelis--Menten \citep{michaelis1913}, or any rate law fixed by
chemistry) whose parameters are, by hypothesis (\Cref{sec:runtime}), fixed by
physical constants rather than fitted.
\item $\Probe$ is the \emph{perturbation space}: admissible inputs (compounds
of known composition and their dosing profiles).
\item $\Rd\colon\Man\to\mathcal{Y}$ is the \emph{readout map} into an
observation space $\mathcal{Y}$: what an instrument reports of a state. The
canonical readout is an abundance distribution over measurable species,
$\Rd(\st)=(\text{abundance of species }s\text{ in state }\st)_{s}$.
\item $\Dom$ is a finite \emph{oscillator index set}: the physiological
subsystems (\Cref{sec:coherence}) whose phase-consistency defines health.
\end{itemize}
\end{definition}

\begin{definition}[Probe; fate; trace]
\label{def:probe}
A \emph{probe} $\probe\in\Probe$ is a perturbation of \emph{known
composition}: a compound (or timed sequence) whose molecular identity and
quantity are exactly specified. Given a state $\st_0$, the \emph{fate} of
$\probe$ is the trajectory $\st(t;\st_0,\probe)$ obtained by integrating
$\dot\st=F(\st,\probe)$ from $\st_0$. The \emph{trace} is the time-indexed
readout $\trace(\st_0,\probe)=(\Rd(\st(t;\st_0,\probe)))_{t\ge0}$; a single
measurement returns the reading $\Rd(\st(t^{\ast}))$ at a read time
$t^{\ast}$.
\end{definition}

\begin{remark}[A probe is a meal of known composition]
\label{rem:meal}
Nothing in \Cref{def:probe} requires a probe to be exotic or inert. Any
administered substance perturbs $\Sigma$; a meal is such a perturbation of
\emph{unknown} composition, hence of unpredictable fate. A probe is a
perturbation of the same physical class whose composition is known, and it is
that knowledge---not inertness---that makes its fate a usable measurement. The
perturbation is not a hazard to be avoided but the instrument itself. No
theorem depends on this reading.
\end{remark}

\subsection{Two structural assumptions}

We adopt two premises and carry them throughout. They are the only assumptions
of the paper about the observer and the state; everything else is derived.

\begin{assumption}[Finite resolution]
\label{ass:finite}
Any physical observer of $\Sigma$ has finite resolving power: there is a
$\delta>0$ such that states within distance $\delta$ produce observationally
indistinguishable readouts,
$\dist(\st,\st')<\delta\Rightarrow\norm{\Rd(\st)-\Rd(\st')}$ below the
instrument's discrimination threshold.
\end{assumption}

\begin{assumption}[Non-exhaustibility]
\label{ass:nonexhaust}
No finite battery of measurements exhausts the state: for every finite set of
readouts there remains a distinction the system could exhibit that the battery
does not resolve. Equivalently, at every scale there is a finer sub-state the
current resolution collapses. This is an order condition on refinement (no
terminal stage), not a claim of infinite dimension.
\end{assumption}

\begin{remark}[Why non-exhaustibility, not infinitude]
\label{rem:why-nonexhaust}
\Cref{ass:nonexhaust} is deliberately an order condition and not a cardinality
claim. A living system is coupled to an environment, a history, and a
regulatory context that no finite instrument battery closes off; there is
always a further assay, a further time point, a further covariate. It is
exactly the impossibility of \emph{finishing} the characterisation of a state
against everything it is not that forces the positive floor of \Cref{sec:floor},
and a completed totality is neither needed nor the right object.
\end{remark}

% =====================================================================
\section{The Resolution Floor}
\label{sec:floor}
% =====================================================================

The first quantitative fact is that resolution has a strictly positive floor,
and that this floor is not a nuisance to be engineered away but the width of
the region within which a state counts as determined. It is the reason a
diagnosis is a region and not a point, and the reason \cf{cfc} returns a
\cf{Cell} rather than a label.

\begin{definition}[Separation and residue]
\label{def:residue}
For a state $\st$, its \emph{separation} $\sigma(\st)$ is the least
observational effort---the discrimination distance in $\mathcal{Y}$ realised
over admissible probes---required to distinguish $\st$ from every other state
of $\Man$. The \emph{residue} $\Rd_{\ast}(\st)\ge0$ is the irreducible part of
this separation that no admissible measurement removes; it is $0$ iff the
distinction of $\st$ from the whole of $\Man$ is completed.
\end{definition}

\begin{theorem}[Floor Theorem]
\label{thm:floor}
Under \Cref{ass:finite,ass:nonexhaust} there is a constant $\flo>0$ such that
$\Rd_{\ast}(\st)\ge\flo$ for every state $\st\in\Man$: no state is resolved to
a point.
\end{theorem}

\begin{proof}
Suppose some state $\st$ had zero residue, $\Rd_{\ast}(\st)=0$. By
\Cref{def:residue} this means the distinction of $\st$ from every other state
of $\Man$ is completed by a finite measurement---every state the system could
exhibit has been brought to bear in separating $\st$ from it. But completing
the comparison of $\st$ against all of $\Man\setminus\set{\st}$ exhausts the
system's distinctions at a finite stage, contradicting non-exhaustibility
(\Cref{ass:nonexhaust}). Hence $\Rd_{\ast}(\st)>0$ for every $\st$. Passing to
the quotient of $\Man$ under indistinguishability (\Cref{ass:finite}, which
identifies states closer than $\delta$), the strictly positive residues are
bounded below by their infimum $\flo:=\inf_{\st}\Rd_{\ast}(\st)>0$: were the
infimum $0$ there would be states of arbitrarily small residue, i.e.\ nearly
completed comparisons, and each in the limit would finish the comparison and
exhaust the system---the same contradiction. Thus $\flo>0$.
\end{proof}

\begin{corollary}[Solved is a cell]
\label{cor:cell}
The finest determination of a state is a \emph{cell} of diameter $\flo$: the
set of states an ideal measurement cannot separate from the true one. ``The
state'' is this cell, not its centre.
\end{corollary}

\begin{definition}[Cell]
\label{def:cell}
For $\st\in\Man$, the \emph{cell} of $\st$ is
$\cell(\st):=\set{\st'\in\Man:\dist(\st,\st')<\flo}$, the ball of radius
$\flo$ under the state metric. A subset $A\subseteq\Man$ is \emph{cell-tight}
if $\diam(A)<\flo$, i.e.\ $A$ lies within a single cell.
\end{definition}

\begin{principle}[The target is the cell, not the point, and not a label]
\label{prin:goal}
Two consequences fix the whole language. \textbf{(a)}~Because the target is a
cell, the diagnostic task is \emph{localization}---shrinking the admissible
region to one cell-tight set---not \emph{recognition against a catalogue}.
\textbf{(b)}~Because a cell is defined intrinsically (by indistinguishability),
it needs no external library of named states to be identified: a diagnosis is
complete without naming a disease. \cf{cfc}'s diagnostic output type is
therefore \cf{Cell}, never a nosological label.
\end{principle}

\begin{remark}[What the Floor Theorem does and does not say]
\label{rem:floor-scope}
\Cref{thm:floor} does not assert that instruments are poor; it asserts that a
genuinely non-exhaustible state cannot, even with an ideal instrument,
terminate at a residue-free point. A ``perfect'' diagnosis---one returning a
point---would require completing a comparison against a state that, by
hypothesis, is never completed. \Cref{sec:diagnosis,sec:cascade} turn this into
an honest stopping rule rather than a caveat.
\end{remark}

% =====================================================================
\section{Coherence: The Geometry of Health}
\label{sec:coherence}
% =====================================================================

Diagnosis needs a notion of \emph{where health is} in $\Man$ that does not
appeal to a catalogue. We supply it through coherence: health is
phase-consistency of the oscillator ensemble, a geometric property of the
state, measured by a scalar in $[0,1]$ and refined into a vector over
oscillator classes.

\subsection{Oscillators and the coherence index}

\begin{definition}[Oscillator; performance metric]
\label{def:oscillator}
An \emph{oscillator} $i\in\Dom$ is a physiological subsystem with a
characteristic frequency and a scalar \emph{performance metric}
$\Pi_i\colon\Man\to\R$ (a rate, amplitude, open probability, period stability,
or the like) read from the state. Each oscillator carries an \emph{optimal}
value $\Pi_i^{\mathrm{opt}}$ (performance at full phase-lock) and a
\emph{degraded} value $\Pi_i^{\mathrm{deg}}$ (performance at zero phase-lock),
with $\Pi_i^{\mathrm{opt}}\neq\Pi_i^{\mathrm{deg}}$.
\end{definition}

\begin{definition}[Coherence index]
\label{def:coherence-index}
The \emph{coherence index} of oscillator $i$ at state $\st$ is the clamped
affine normalisation
\begin{equation}
\ecoh_i(\st)\;=\;\operatorname{clip}_{[0,1]}\!
\left(\frac{\Pi_i(\st)-\Pi_i^{\mathrm{deg}}}
           {\Pi_i^{\mathrm{opt}}-\Pi_i^{\mathrm{deg}}}\right),
\qquad
\operatorname{clip}_{[0,1]}(z)=\max(0,\min(1,z)),
\label{eq:coherence}
\end{equation}
so $\ecoh_i(\st)=1$ at optimal performance and $\ecoh_i(\st)=0$ at degraded
performance, monotone in $\Pi_i$ between them.
\end{definition}

\begin{proposition}[Uniqueness of the coherence normalisation]
\label{prop:coherence-unique}
Among affine maps $\R\to[0,1]$ carrying $\Pi_i^{\mathrm{deg}}\mapsto0$ and
$\Pi_i^{\mathrm{opt}}\mapsto1$, \eqref{eq:coherence} is the unique one, and its
clamp is the unique extension to all of $\R$ that is monotone, idempotent on
$[0,1]$, and constant outside.
\end{proposition}

\begin{proof}
An affine map $z\mapsto az+b$ with two prescribed value pairs has $a,b$
determined uniquely by solving the two linear equations, giving the ratio in
\eqref{eq:coherence}. Extending to $\R$ while keeping the image in $[0,1]$,
monotone, and equal to the affine map on the pre-image of $[0,1]$ forces the
value $0$ below and $1$ above, which is exactly $\operatorname{clip}_{[0,1]}$.
\end{proof}

\begin{definition}[Weighted cellular coherence]
\label{def:cell-coherence}
Given nonnegative weights $\set{w_i}_{i\in\Dom}$ with $W=\sum_i w_i>0$, the
\emph{cellular coherence} at state $\st$ is
\begin{equation}
\ecoh(\st)\;=\;\frac{1}{W}\sum_{i\in\Dom} w_i\,\ecoh_i(\st)\;\in[0,1].
\label{eq:cell-coherence}
\end{equation}
The \emph{disease index} of oscillator $i$ is $D_i(\st)=1-\ecoh_i(\st)$, and
the \emph{disease vector} is $\dvec(\st)=(D_i(\st))_{i\in\Dom}\in[0,1]^{|\Dom|}$.
\end{definition}

\begin{proposition}[Coherence bounds and severity]
\label{prop:coherence-bounds}
$\ecoh(\st)\in[0,1]$ for every $\st$, and the total severity
$1-\ecoh(\st)=\tfrac1W\sum_i w_i D_i(\st)$ is the weighted mean disease index.
Both are computable in $\mathcal{O}(|\Dom|)$ from the metrics
$\set{\Pi_i(\st)}$.
\end{proposition}

\begin{proof}
Each $\ecoh_i(\st)\in[0,1]$ by the clamp; a weighted average of numbers in
$[0,1]$ with positive weights lies in $[0,1]$. The severity identity is
immediate from $D_i=1-\ecoh_i$ and linearity. Each $\ecoh_i$ is
$\mathcal{O}(1)$; summing over $\Dom$ is $\mathcal{O}(|\Dom|)$.
\end{proof}

\subsection{The coherent set}

\begin{definition}[Coherent set; health]
\label{def:coherent}
Fix a coherence threshold $\theta\in(0,1]$. The \emph{coherent set} is
\begin{equation}
\Coh_\theta\;=\;\set{\st\in\Man:\ecoh_i(\st)\ge\theta\ \text{for all }i\in\Dom},
\label{eq:coherent}
\end{equation}
the states in which every oscillator is within tolerance of phase-lock. A
state is \emph{healthy} iff it lies in $\Coh_\theta$. Where the threshold is
fixed we write $\Coh$.
\end{definition}

\begin{remark}[Health is a region, not a catalogue complement]
\label{rem:health-region}
$\Coh$ is defined \emph{intrinsically}, by an inequality on the state's own
oscillator metrics, with no reference to any list of diseases. Disease is then
simply $\Man\setminus\Coh$: the states failing coherence. This is the geometric
content of ``disease is loss of phase-consistency,'' and it is what lets
diagnosis proceed by localization against $\Coh$ rather than by matching
against named entries. No theorem depends on any clinical reading of the word
``healthy.''
\end{remark}

\begin{lemma}[Coherence is Lipschitz where the metrics are]
\label{lem:coherence-lipschitz}
If each $\Pi_i$ is $L_i$-Lipschitz on $\Man$, then $\ecoh$ is
$L$-Lipschitz with
$L=\tfrac1W\sum_i w_i L_i/\abs{\Pi_i^{\mathrm{opt}}-\Pi_i^{\mathrm{deg}}}$, and
each cell $\cell(\st)$ has coherence variation at most $L\flo$.
\end{lemma}

\begin{proof}
The clamp is $1$-Lipschitz and the affine map has slope
$1/\abs{\Pi_i^{\mathrm{opt}}-\Pi_i^{\mathrm{deg}}}$, so
$\ecoh_i$ is $(L_i/\abs{\Pi_i^{\mathrm{opt}}-\Pi_i^{\mathrm{deg}}})$-Lipschitz;
the weighted average is $L$-Lipschitz by the triangle inequality. Over a cell
of diameter $\flo$ the coherence changes by at most $L\flo$.
\end{proof}

% =====================================================================
\section{Diagnosis as Localization}
\label{sec:diagnosis}
% =====================================================================

We now make precise the sense in which a probe's reading \emph{locates} the
patient, and the exact condition under which a probe is diagnostically
sufficient. Diagnosis is the shrinking of an admissible region to one cell; it
succeeds when the map from state to reading loses no information above the
floor.

\begin{definition}[State-resolved forward model]
\label{def:stateresolved}
A forward model of $\Sigma$ is \emph{state-resolved} if it predicts the trace
$\trace(\st_0,\probe)$ as a function of the initial state $\st_0$ over all of
$\Man$, not merely at a reference $\st^{\mathrm{ref}}$. Write
$\fwd_\probe\colon\Man\to\mathcal{Y}$,
$\fwd_\probe(\st_0)=\Rd(\st(t^{\ast};\st_0,\probe))$, for the predicted reading
map of probe $\probe$ at read time $t^{\ast}$.
\end{definition}

\begin{remark}[Why a reference-only model is insufficient]
\label{rem:reference}
A model valid only at $\st^{\mathrm{ref}}$ predicts one reading; a measured
deviation certifies only \emph{that} the patient differs, not \emph{how}.
Inversion---reading the state off the reading---requires the reading to vary
with the state, i.e.\ $\fwd_\probe$ as a genuine map of $\st_0$. This is the
crux of localization and the sole substantive demand the method places on the
model (discharged in \Cref{sec:runtime}).
\end{remark}

\begin{definition}[Localization map]
\label{def:loc}
Given probe $\probe$ and a measured reading $y\in\mathcal{Y}$ at time
$t^{\ast}$, the \emph{localization} induced by $\probe$ is the preimage
\begin{equation}
\loc_\probe(y)\;=\;\set{\st_0\in\Man:\norm{\fwd_\probe(\st_0)-y}<\delta},
\label{eq:loc}
\end{equation}
the set of states whose predicted reading matches the measurement to within the
resolution $\delta$ (\Cref{ass:finite}). Reading the probe replaces the prior
admissible region $A$ by $A\cap\loc_\probe(y)$.
\end{definition}

\begin{theorem}[Localization Theorem / diagnostic sufficiency]
\label{thm:localization}
A probe $\probe$ is \emph{diagnostically sufficient} for a region
$A\subseteq\Man$ if and only if its reading map restricted to $A$ is injective
down to the floor: for all $\st_0,\st_0'\in A$,
\begin{equation}
\norm{\fwd_\probe(\st_0)-\fwd_\probe(\st_0')}<\delta
\;\Longrightarrow\;\dist(\st_0,\st_0')<\flo.
\label{eq:floor-injective}
\end{equation}
Under this condition, for every measurement $y$ the localization
$\loc_\probe(y)\cap A$ is cell-tight, so the probe pins the state to one cell.
Conversely, if \eqref{eq:floor-injective} fails, there are two states at
distance $\ge\flo$ producing the same reading within $\delta$, and no single
reading of $\probe$ separates their cells.
\end{theorem}

\begin{proof}
($\Leftarrow$) Assume \eqref{eq:floor-injective}. Take
$\st_0,\st_0'\in\loc_\probe(y)\cap A$. Then
$\norm{\fwd_\probe(\st_0)-y}<\delta$ and $\norm{\fwd_\probe(\st_0')-y}<\delta$,
so $\norm{\fwd_\probe(\st_0)-\fwd_\probe(\st_0')}<2\delta$; applying
\eqref{eq:floor-injective} with the instrument threshold set at $2\delta$ (or,
equivalently, absorbing the factor into $\delta$, as \Cref{ass:finite}
tolerates a single fixed discrimination scale) yields
$\dist(\st_0,\st_0')<\flo$. Hence any two states in the localization are within
$\flo$, so $\diam(\loc_\probe(y)\cap A)<\flo$: the set is cell-tight.
($\Rightarrow$) Suppose \eqref{eq:floor-injective} fails: there exist
$\st_0,\st_0'\in A$ with
$\norm{\fwd_\probe(\st_0)-\fwd_\probe(\st_0')}<\delta$ yet
$\dist(\st_0,\st_0')\ge\flo$. Set $y=\fwd_\probe(\st_0)$. Then both states lie
in $\loc_\probe(y)$ (the second because its reading is within $\delta$ of $y$),
so the localization contains two points at distance $\ge\flo$, meeting two
distinct cells; no single reading distinguishes them, and $\probe$ is
insufficient for $A$.
\end{proof}

\begin{corollary}[Probe design is an injectivity problem]
\label{cor:design}
Designing a diagnostic probe (or panel, \Cref{sec:cascade}) for a region $A$ is
choosing a probe whose reading map is floor-injective on $A$. Diagnosis is then
evaluation of $\loc_\probe$ at the measured reading. No disease library enters;
the output is the cell (\Cref{prin:goal}).
\end{corollary}

\begin{corollary}[Coherence is read from the cell]
\label{cor:coherence-from-cell}
If $\probe$ is diagnostically sufficient for $A$ and each metric $\Pi_i$ is
$L_i$-Lipschitz, then reading $\probe$ determines every coherence index
$\ecoh_i$ and hence the disease vector $\dvec$ to within $L_i\flo$
(\Cref{lem:coherence-lipschitz}): the localized cell fixes the health geometry
up to the floor.
\end{corollary}

\begin{proof}
By \Cref{thm:localization} the localization is cell-tight, of diameter $<\flo$;
by \Cref{lem:coherence-lipschitz} each $\ecoh_i$ varies by at most $L_i\flo$
over a cell, so its value is fixed to that tolerance, and $\dvec$ with it.
\end{proof}

% =====================================================================
\section{Cascading Probes and the Trajectory Certificate}
\label{sec:cascade}
% =====================================================================

A single reading of a single probe may fail to be floor-injective; more
information is needed. Rather than administer many independent probes, we
exploit that each physiological product is itself a reactive species, so a
probe can generate its own successors---a \emph{cascade}---whose structure
yields both richer localization and a strong non-forgeability property that a
mere agreement of magnitudes cannot provide.

\begin{definition}[Cascade probe; panel]
\label{def:cascade}
A \emph{cascade probe} is a probe $\probe$ whose fate passes through a designed
sequence of species $\probe=P_0\to P_1\to P_2\to\cdots$, where each $P_{k+1}$
is a product of a reaction consuming $P_k$, and each transition $P_k\to
P_{k+1}$ is \emph{state-gated}: whether $P_{k+1}$ forms in observable quantity
within the read window depends on the state $\st$. A \emph{panel} is a finite
family of probes (cascade or single) read together.
\end{definition}

\begin{definition}[Traversed path; committed record]
\label{def:trajectory}
The \emph{traversed path} of a cascade in state $\st_0$ is the maximal prefix
$P_0\to\cdots\to P_{d(\st_0)}$ of state-gated transitions proceeding to
observable extent; $d(\st_0)$ is the \emph{depth}. Because each transition
consumes its predecessor and produces its successor, the species population is
a \emph{monotone committed record}: a later species' presence is not undone by
subsequent dynamics on the read timescale, so the reading reports the path, not
only the instantaneous state.
\end{definition}

\begin{theorem}[Certificate Theorem]
\label{thm:certificate}
Let $P_k$ be a species that forms in observable quantity only if every
transition $P_0\to\cdots\to P_k$ has occurred. Then observation of $P_k$
\emph{entails} that the state $\st_0$ enables the whole path: the set of states
consistent with the reading is contained in
\begin{equation}
\bigcap_{j\le k}\set{\st:\text{transition }j\text{ is state-enabled at }\st}.
\label{eq:consistent-set}
\end{equation}
Consequently, a hypothesised state $\hat\st$ that fails to enable some
transition $j\le k$ is \emph{refuted} by the presence of $P_k$: no such
$\hat\st$ can reproduce the reading.
\end{theorem}

\begin{proof}
The forward implication is the existence condition of \Cref{def:cascade}:
$P_k$ observable $\Rightarrow$ each prior transition proceeded $\Rightarrow$
each prior transition was state-enabled at $\st_0$. Hence the consistent-state
set is the intersection of the enabling sets of transitions $1,\dots,k$, which
is \eqref{eq:consistent-set}. If $\hat\st$ lies outside the enabling set of
some $j\le k$, then at $\hat\st$ transition $j$ does not proceed, so $P_j$---and
therefore $P_k$, which requires $P_j$---is absent; observing $P_k$ contradicts
$\hat\st$.
\end{proof}

\begin{principle}[Existence beats consistency]
\label{prin:existence}
\Cref{thm:certificate} is a stronger epistemic guarantee than agreement of
magnitudes. A consistency check asks whether a hypothesis \emph{matches}
observed values; it can be satisfied by a well-chosen falsehood. A cascade
reading asks whether the observed species \emph{could have formed}; a falsehood
failing to enable a transition does not merely mismatch a number---it predicts
a species that is \emph{absent}, and absence cannot be argued around. The
terminal products are a constructive proof of their own prerequisites.
\end{principle}

\begin{theorem}[Rigidity: the correct halting criterion]
\label{thm:rigidity}
Let a panel yield observed species $S=\set{P_k}$, and define the
\emph{surviving set}
\begin{equation}
\Man_S\;=\;\bigcap_{P_k\in S}\set{\st:\text{the path to }P_k\text{ is enabled at }\st}.
\label{eq:surviving}
\end{equation}
Then $\Man_S$ is non-increasing as $S$ grows, and the state is
\emph{sufficiently localized} exactly when $\Man_S$ is cell-tight
($\diam(\Man_S)<\flo$): every account of the state outside a single cell is
refuted by the absence it would predict of some observed species. This is a
\emph{rigidity} condition---no alternative account survives---not a threshold on
any scalar.
\end{theorem}

\begin{proof}
Each observed species adds a constraint $\set{\st:\text{its path enabled}}$;
$\Man_S$ is the intersection of these constraints, and an intersection is
non-increasing under the addition of sets, so $\Man_S$ shrinks (weakly) as $S$
grows. When $\diam(\Man_S)<\flo$, $\Man_S$ lies within one cell
(\Cref{cor:cell}); any $\st$ outside that cell violates at least one constraint,
i.e.\ predicts some observed $P_k$ absent, and is refuted by
\Cref{thm:certificate}. Conversely, if $\diam(\Man_S)\ge\flo$, two states in
$\Man_S$ occupy distinct cells and both reproduce the reading, so the state is
not localized. Sufficiency is therefore exactly cell-tightness.
\end{proof}

\begin{corollary}[Single-measurement richness]
\label{cor:single}
A readout returning the full species distribution---the whole population of
cascade products present, not one terminal product---delivers the entire
constraint family $S$ in one measurement, hence $\Man_S$ in one shot. The
instrument's report \emph{is} the committed record of the traversed path.
\end{corollary}

\begin{remark}[A continuously working probe]
\label{rem:continuous}
Because a cascade is self-propagating and its record monotone, a single
administration installs a process that keeps advancing and recording after
administration. The state is therefore readable at any later time by any
observer with the instrument: the probe does not deliver a value once and
expire; it remains a live, queryable source until the cascade completes.
``Continuous'' is here a property of the probe---always working---not of the
measurement schedule. No theorem depends on this reading.
\end{remark}

% =====================================================================
\section{Therapy as Gap-Closure Without a Named Disease}
\label{sec:therapy}
% =====================================================================

We now define treatment purely in terms of the localized state and coherence,
with no reference to a disease label. Therapy closes the measured gap to the
coherent set by the sparsest characterised interventions, executed as a
contraction.

\begin{definition}[Gap]
\label{def:gap}
For a localized state $\st$, the \emph{gap} is the displacement to the nearest
coherent state,
\begin{equation}
\gap(\st)\;=\;\coh-\st,\qquad\coh=\argmin_{z\in\Coh}\dist(\st,z),
\label{eq:gap}
\end{equation}
a measured vector defined without naming what caused the deviation. (If the
minimiser is non-unique, fix any measurable selection; no theorem depends on
the choice.)
\end{definition}

\begin{definition}[Intervention basis]
\label{def:basis}
An \emph{intervention basis} $\basis=\set{\intv_1,\dots,\intv_m}$ is a
collection of characterised compounds, each with a known state-resolved effect:
applying $\intv_i$ at state $\st$ produces a local displacement
$e_i(\st)\in\R^n$, the compound's effect at the current state computed by the
forward model. Compounds are indexed by their displacement vectors
(\emph{properties}), not by indication.
\end{definition}

\begin{remark}[Repurposing is the default, not the exception]
\label{rem:repurpose}
Because a compound enters therapy through its displacement $e_i(\st)$ and not
through any indication, any compound whose displacement has a component along
the gap is admissible for any patient. The distinction
``on-label / off-label'' does not arise: there are no labels, only properties
and gaps, and the pharmacopeia is addressable by its effects
\citep{pushpakom2019drug}. No theorem depends on this reading.
\end{remark}

\begin{problem}[Minimal gap-closing intervention]
\label{prob:closure}
Given a localized state $\st$ with gap $\gap(\st)$ and basis $\basis$, choose
nonnegative doses $\dose\in\Rp^{m}$ to close the gap with the fewest compounds:
\begin{equation}
\min_{\dose\in\Rp^{m}}\ \norm{\dose}_1
\quad\text{s.t.}\quad
\Bigl\lVert\gap(\st)-\textstyle\sum_i \dose_i\,e_i(\st)\Bigr\rVert\le\flo.
\label{eq:closure}
\end{equation}
The $\ell_1$ objective selects a sparse intervention set; the constraint
requires the combined local displacement to close the gap to within the floor.
\end{problem}

\begin{proposition}[Well-posedness and sparsity]
\label{prop:closure-wellposed}
\Cref{prob:closure} is a convex program; it is feasible iff $\gap(\st)$ lies
within $\flo$ of the cone $\operatorname{cone}\set{e_i(\st)}$, and when feasible
it attains a minimiser. If the vectors $\set{e_i(\st)}$ are in general position,
a minimiser is supported on at most $n$ compounds.
\end{proposition}

\begin{proof}
The objective $\norm{\dose}_1$ is convex and the constraint set
$\set{\dose\ge0:\norm{\gap-\sum_i \dose_i e_i}\le\flo}$ is the intersection of
the nonnegative orthant with the preimage of a closed ball under an affine map,
hence closed and convex; the problem is convex
\citep{boyd2004convex}. Feasibility is exactly the stated cone condition by
definition of the constraint. On the feasible set the objective is coercive on
the orthant (grows with $\norm{\dose}_1$) and continuous, and the feasible set
is closed, so a minimiser is attained. Sparsity: the feasible set is a
polyhedron intersected with a ball; by the representation of basic feasible
solutions of the linear relaxation, and the sparsity-inducing geometry of the
$\ell_1$ ball \citep{tibshirani1996lasso,candes2006robust}, a minimiser lies on
a face of dimension $\le n$, so at most $n$ doses are nonzero when the
$e_i(\st)$ are in general position.
\end{proof}

The local displacements $e_i(\st)$ are exact only at $\st$; applying the doses
moves the state, after which the displacements change. Therapy is therefore
iterative.

\begin{definition}[Therapeutic map; contraction]
\label{def:contraction}
Let $T\colon\Man\to\Man$ send a state to the state resulting from localizing
it (\Cref{sec:diagnosis}), solving \Cref{prob:closure}, and applying the chosen
intervention. $T$ is a \emph{gap contraction} on a neighbourhood $U$ of $\Coh$
if there is $\lambda\in[0,1)$ with $\dist(T\st,\Coh)\le\lambda\,\dist(\st,\Coh)$
for all $\st\in U$.
\end{definition}

\begin{theorem}[Convergence to coherence]
\label{thm:convergence}
If $T$ is a gap contraction on $U$ with modulus $\lambda$, then iterating $T$
from any $\st_0\in U$ drives the state to the coherent set: the sequence
$\st_{k+1}=T\st_k$ satisfies
$\dist(\st_k,\Coh)\le\lambda^{k}\dist(\st_0,\Coh)\to0$, and reaches
$\dist(\st_k,\Coh)\le\flo$ after at most
$k^{\ast}=\bigl\lceil\log_{1/\lambda}(\dist(\st_0,\Coh)/\flo)\bigr\rceil$
steps, at which point the floor forbids further resolvable contraction. No
global linearity of drug effects is required---only that each measured,
locally-composed step contracts the gap.
\end{theorem}

\begin{proof}
By the contraction inequality and induction,
$\dist(\st_k,\Coh)\le\lambda^{k}\dist(\st_0,\Coh)$; since $\lambda\in[0,1)$ this
tends to $0$. Solving $\lambda^{k}\dist(\st_0,\Coh)\le\flo$ for $k$ gives
$k\ge\log_{1/\lambda}(\dist(\st_0,\Coh)/\flo)$, i.e.\ the stated $k^{\ast}$.
Below $\flo$ the states are indistinguishable (\Cref{ass:finite}), so no
further contraction is resolvable and the iteration halts. This is the standard
fixed-point iteration of a contraction \citep{banach1922}, specialised to
distance-to-a-set with the floor as halting threshold.
\end{proof}

\begin{principle}[Cure without naming]
\label{prin:cure}
\Cref{thm:convergence} restores coherence---closes the gap---without ever
identifying a disease. The therapeutic objective is the measured residual gap,
and success is its reduction to the floor. Naming, if desired, is a downstream
description of which cell the state occupied; it is never a prerequisite for
treatment.
\end{principle}

% =====================================================================
\section{Diagnosis and Therapy Are One Operation}
\label{sec:unity}
% =====================================================================

The measurement primitive and the therapeutic primitive are the same object
read in two directions. This is what lets \cf{cfc} expose a single verb for
both and makes the treatment loop self-monitoring.

\begin{theorem}[Measurement--action identity]
\label{thm:unity}
Every characterised perturbation $\intv\in\Probe$ acts simultaneously as a
measurement and as an intervention:
\begin{enumerate}[label=\textup{(\roman*)},leftmargin=2.2em,itemsep=2pt]
\item as a \emph{measurement}, its trace $\trace(\st,\intv)$ localizes the
state (\Cref{thm:localization});
\item as an \emph{intervention}, its displacement moves the state
(\Cref{def:basis});
\end{enumerate}
and these are not two applications but one: applying $\intv$ yields a trace
(information) \emph{by} moving the state (action). A probe is a compound read
for its fate; a drug is a compound applied for its displacement; they are the
same primitive, and each does both.
\end{theorem}

\begin{proof}
A perturbation induces a single trajectory (\Cref{def:probe}). Its readout is
the measurement (i); its endpoint is the action (ii). Both are functions of the
one applied $\intv$: the map $\intv\mapsto(\text{trace},\text{displacement})$
has one input and two outputs computed from the same induced trajectory.
\end{proof}

\begin{corollary}[Self-terminating therapeutic loop]
\label{cor:loop}
Because the therapeutic step is also a measurement, the contraction of
\Cref{sec:therapy} needs no external monitor: each applied intervention
re-localizes the state as a side effect of acting. The loop
measure--act--remeasure is a single iterated perturbation, and it halts by the
rigidity criterion (\Cref{thm:rigidity})---when the accumulated record
over-determines the state to within the floor---which is simultaneously the
diagnostic sufficiency condition and the therapeutic success condition.
\end{corollary}

\begin{theorem}[Wellness Theorem]
\label{thm:wellness}
Combining \Cref{thm:rigidity,thm:convergence}, the terminal condition of the
\cf{cfc} loop is a state that is both \emph{coherent} (its gap is below the
floor) and \emph{rigidly known} (no alternative account survives its own
probes). This double condition---and only it---is the operative definition of
wellness in the calculus: a well system returns its probes at full resolution
and admits no gap-closing intervention beyond the floor.
\end{theorem}

\begin{proof}
Termination of the loop requires both sub-conditions. If the gap exceeds the
floor, \Cref{thm:convergence} supplies a further contracting step, so the loop
continues; hence at termination $\dist(\st,\Coh)\le\flo$ (coherent). If more
than one cell survives, \Cref{thm:rigidity} shows a further probe can refute a
rival account, so the loop continues; hence at termination the surviving set is
cell-tight (rigidly known). Conversely, when both hold, no probe changes the
localized cell and no intervention resolvably reduces the gap, so the loop is at
a fixed point. The two conditions are thus jointly necessary and sufficient for
termination.
\end{proof}

\begin{remark}[Concern as the negation of wellness]
\label{rem:concern}
The language's name is the negation of \Cref{thm:wellness}: a
\emph{cause for concern} is any surviving cell whose coherence falls below
threshold, i.e.\ any $\st$ in the localized set with $\ecoh(\st)<\theta$. A
\cf{cfc} program's \cf{concern} value carries exactly this---the localized
cell together with its coherence deficit and disease vector---never a disease
name. No theorem depends on the reading.
\end{remark}

% =====================================================================
\section{The Runtime Must Be Derived, Not Fitted}
\label{sec:runtime}
% =====================================================================

Every construction above presupposes a state-resolved forward model
(\Cref{def:stateresolved}). We now state the one condition this model must
meet, as a theorem about extrapolation, because the safety of \emph{in silico}
experimentation---the whole point of running a therapy search against a model
before touching a patient---reduces to it.

\begin{definition}[Fitted vs.\ derived model]
\label{def:fitvsder}
A forward model is \emph{fitted} if its rate parameters are estimated to
minimise error against observations drawn from a training region
$A_{\mathrm{tr}}\subseteq\Man$. It is \emph{derived} if its rate laws and their
parameters are fixed by physical law (thermodynamic and kinetic constants,
conservation constraints) independently of any observation region
\citep{palsson2015systems,orth2010fba}.
\end{definition}

\begin{theorem}[Extrapolation gap of fitted models]
\label{thm:extrap}
A fitted model carries no guarantee of accuracy outside $A_{\mathrm{tr}}$:
there exist state systems agreeing with the training data on $A_{\mathrm{tr}}$
to any tolerance yet disagreeing arbitrarily on $\Man\setminus A_{\mathrm{tr}}$.
Consequently a fitted model's reading map $\fwd_\probe$ is trustworthy only
near $A_{\mathrm{tr}}$, whereas a derived model, whose $F$ is fixed pointwise on
$\Man$ by law, admits a uniform error bound over $\Man$.
\end{theorem}

\begin{proof}
Standard non-identifiability. The training data constrain $F$ only on
$A_{\mathrm{tr}}$: any extension of $F$ agreeing there is admissible, and two
admissible extensions may be chosen to diverge by any prescribed amount off
$A_{\mathrm{tr}}$ (add to one a smooth bump supported in
$\Man\setminus A_{\mathrm{tr}}$ of arbitrary height). Hence the fitted
$\fwd_\probe$, obtained by integrating $F$, is pinned only near
$A_{\mathrm{tr}}$. A derived $F$ is not an extension chosen to fit but a
function specified on all of $\Man$ by its physical form, so its accuracy is
governed by the fidelity of the physical law, uniformly, not by proximity to
data.
\end{proof}

\begin{principle}[Free exploration requires a derived runtime]
\label{prin:derived}
The value of \cf{cfc} is that a clinician may experiment freely against the
model---searching the intervention basis, testing cascades, exploring the
gap---at no risk to the patient, applying to the patient only the intervention
the search selects. But free exploration \emph{is} off-training-region
evaluation by definition. By \Cref{thm:extrap} a fitted runtime becomes
unreliable exactly where the search explores; only a derived runtime supports
risk-free \emph{in silico} search. The requirement is not stylistic: the safety
of the method reduces to the extrapolation validity of the runtime. A \cf{cfc}
runtime is \emph{sound} only if its forward model is derived
(\Cref{def:soundness}).
\end{principle}

% =====================================================================
\section{Lexical Structure}
\label{sec:lexical}
% =====================================================================

A \cf{cfc} source file is UTF-8 text with extension \cf{.cfc}. Lexing produces
a token stream from the classes below; tokenisation follows maximal munch, with
keywords taking priority over identifiers \citep{hopcroft2006automata}.

\begin{definition}[Token classes]
\label{def:tokens}
\leavevmode
\begin{itemize}[leftmargin=1.4em,itemsep=1pt]
\item \textbf{Comments.} \cf{--} to end of line; discarded, lexically
equivalent to whitespace.
\item \textbf{Keywords.} \cf{floor}, \cf{patient}, \cf{probe}, \cf{panel},
\cf{measure}, \cf{localize}, \cf{cell}, \cf{gap}, \cf{close}, \cf{toward},
\cf{coherent}, \cf{using}, \cf{basis}, \cf{until}, \cf{resolved},
\cf{contested}, \cf{decline}, \cf{concern}, \cf{when}, \cf{emit}, \cf{assert},
\cf{yield}, \cf{let}, \cf{in}, \cf{as}, \cf{with}, \cf{import}, \cf{module},
\cf{export}, \cf{monitor}, \cf{every}, \cf{observe}, \cf{against},
\cf{reference}.
\item \textbf{Identifiers.} \cf{[A-Za-z\_][A-Za-z0-9\_]*}, excluding keywords.
\item \textbf{Numbers.} decimal integers and floats with optional exponent and
optional floor annotation (\Cref{not:floor-literal}).
\item \textbf{Strings.} double-quoted, with the usual escapes.
\item \textbf{Operators and punctuation.} \cf{:=} (bind), \cf{>>} (cascade,
sequential composition of probe stages), \cf{||} (parallel panel dispatch),
\cf{( ) \{ \} [ ] , : . > < >= <= ==} and the field-access operator \cf{.}\ .
\end{itemize}
Whitespace is insignificant except as a token separator; \cf{cfc} is not
layout-sensitive.
\end{definition}

\begin{notation}[Floor-annotated literals]
\label{not:floor-literal}
A numeric literal may carry an explicit floor annotation with the suffix form
\cf{value\#floor} (read ``value at floor''); e.g.\ \cf{0.82\#1e-3}. A literal
without an explicit floor inherits the ambient program floor set by the most
recent \cf{floor} declaration in scope. There is no way to write a literal of
zero floor (\Cref{thm:no-zero}).
\end{notation}

% =====================================================================
\section{Context-Free Grammar}
\label{sec:grammar}
% =====================================================================

We give the grammar in extended Backus--Naur form
\citep{backus1959,naur1963,aho2006compilers}. Nonterminals are
\textit{italicised}; terminals are \cf{typewriter}; \cf{\{x\}} is zero-or-more
repetition, \cf{[x]} optional, \cf{|} alternation. The start symbol is
\textit{program}.

\begin{definition}[\cf{cfc} grammar]
\label{def:grammar}
\begin{align*}
\textit{program} &\bnf \{\textit{decl}\}\\[2pt]
\textit{decl} &\bnf \textit{floorDecl}\alt\textit{import}\alt\textit{moduleDecl}
   \alt\textit{basisDecl}\alt\textit{patientDecl}\\[2pt]
\textit{floorDecl} &\bnf \cf{floor}\ \textit{number}\\
\textit{import} &\bnf \cf{import}\ \textit{qname}\\
\textit{moduleDecl} &\bnf \cf{module}\ \textit{ident}\ \cf{\{}\ \{\textit{decl}\}\ \cf{\}}\\[2pt]
\textit{basisDecl} &\bnf \cf{basis}\ \textit{ident}\ \cf{\{}\ \{\textit{probeDef}\}\ \cf{\}}\\
\textit{probeDef} &\bnf \cf{probe}\ \textit{ident}\ \cf{\{}\ \textit{spec}\ \cf{\}}\\
\textit{spec} &\bnf \cf{reference}\cf{:}\ \textit{expr}\ \cf{observe}\cf{:}\ \textit{stageChain}\\[2pt]
\textit{patientDecl} &\bnf \cf{patient}\ \textit{ident}\ \cf{\{}\ \{\textit{stmt}\}\ \cf{\}}\\[2pt]
\textit{stmt} &\bnf \textit{bind}\alt\textit{measureStmt}\alt\textit{diagnoseStmt}
   \alt\textit{treatStmt}\\
  &\alt\textit{monitorStmt}\alt\textit{assertStmt}\alt\textit{expr}\\[2pt]
\textit{bind} &\bnf [\cf{let}]\ \textit{ident}\ \cf{:=}\ \textit{expr}\\[2pt]
\textit{measureStmt} &\bnf \cf{measure}\ \textit{panelRef}\ \cf{yield}\ \textit{ident}\\[2pt]
\textit{diagnoseStmt} &\bnf \cf{localize}\ \textit{ident}\\
  &\quad \cf{using}\ \textit{panelRef}\\
  &\quad \cf{until}\ \textit{admit}\\
  &\quad \cf{yield}\ \textit{ident}\\[2pt]
\textit{treatStmt} &\bnf \cf{close}\ \cf{gap}\ \cf{of}\ \textit{ident}\\
  &\quad \cf{toward}\ \cf{coherent}\\
  &\quad \cf{using}\ \cf{basis}\ \textit{ident}\\
  &\quad \cf{until}\ \textit{admit}\\
  &\quad \cf{yield}\ \textit{ident}\\[2pt]
\textit{monitorStmt} &\bnf \cf{monitor}\ \textit{ident}\ \cf{every}\ \textit{number}\ \cf{\{}\ \{\textit{stmt}\}\ \cf{\}}\\[2pt]
\textit{admit} &\bnf \cf{resolved}\ [\cf{otherwise}\ \cf{decline}]\\
  &\alt\cf{coherent}\ [\cf{otherwise}\ \cf{decline}]\\
  &\alt\textit{cond}\\[2pt]
\textit{panelRef} &\bnf \cf{panel}\ \cf{\{}\ \textit{stageChain}\ \cf{\}}\alt\textit{qname}\\
\textit{stageChain} &\bnf \textit{probeRef}\ \{\cf{>>}\ \textit{probeRef}\}\\
  &\alt\textit{stageChain}\ \cf{||}\ \textit{stageChain}\\
\textit{probeRef} &\bnf \textit{qname}\ \cf{(}\ [\textit{argList}]\ \cf{)}\\[2pt]
\textit{assertStmt} &\bnf \cf{assert}\ \textit{cond}\ [\cf{emit}\ \textit{string}]\\[2pt]
\textit{cond} &\bnf \textit{expr}\ \textit{relop}\ \textit{expr}\\
\textit{relop} &\bnf \cf{>}\alt\cf{<}\alt\cf{>=}\alt\cf{<=}\alt\cf{==}\\
\textit{argList} &\bnf \textit{arg}\ \{\cf{,}\ \textit{arg}\}\\
\textit{arg} &\bnf [\textit{ident}\ \cf{:}]\ \textit{expr}\\
\textit{qname} &\bnf \textit{ident}\ \{\cf{.}\ \textit{ident}\}\\
\textit{expr} &\bnf \textit{ident}\alt\textit{number}\alt\textit{string}
   \alt\textit{probeRef}\alt\textit{expr}\cf{.}\textit{ident}\alt\cf{(}\ \textit{expr}\ \cf{)}\\
\textit{type} &\bnf \cf{Cell}\alt\cf{Coherence}\alt\cf{DiseaseVector}\alt\cf{Gap}
   \alt\cf{Probe}\alt\cf{Panel}\\
  &\alt\cf{Intervention}\alt\cf{Reading}\alt\cf{Verdict}\alt\cf{Concern}
\end{align*}
\end{definition}

\begin{remark}[Reading the core forms]
\label{rem:core-forms}
A \cf{cfc} program has three verbs, one per operational theorem.
\cf{measure} administers a panel and binds its \cf{Reading}
(\Cref{def:probe}). \cf{localize \ldots until resolved} runs the rigidity loop
of \Cref{thm:rigidity}, binding a \cf{Cell}; its optional
\cf{otherwise decline} names the contested-closure outcome
(\Cref{thm:decline}). \cf{close gap \ldots toward coherent until coherent}
runs the gap-closing contraction of \Cref{thm:convergence}, binding an
\cf{Intervention}; \cf{coherent} means ``gap below the floor,'' not a numeric
threshold. \cf{monitor \ldots every $\Delta t$} re-runs its body on a schedule,
realising the continuously working probe of \Cref{rem:continuous}.
\cf{assert} is a pre-flight check that may halt with a message.
A \cf{patient} block is the unit of orchestration: a sequence of these
statements whose yields may feed later ones.
\end{remark}

\begin{theorem}[Unambiguity of the core grammar]
\label{thm:unambiguous}
The grammar of \Cref{def:grammar} is unambiguous, and is $LL(1)$ after
left-factoring the \textit{stmt} and \textit{admit} alternatives, each of which
is selected by a distinct leading keyword.
\end{theorem}

\begin{proof}
Every statement form begins with a distinct keyword (\cf{measure},
\cf{localize}, \cf{close}, \cf{monitor}, \cf{assert}, \cf{let}) or is an
\textit{expr}; the leading token thus selects the production, giving $LL(1)$
selection with one token of lookahead. The \textit{admit} alternatives are
distinguished by their leading keyword (\cf{resolved}, \cf{coherent}, or an
\textit{expr} beginning a \textit{cond}). The optional
\cf{otherwise decline} is bound to its enclosing \textit{admit} by the nearest
enclosing rule, with no dangling ambiguity because each \textit{admit} occurs
in exactly one statement. The expression grammar is a standard
member/call-suffix cascade, unambiguous by the usual argument
\citep{aho2006compilers}. Hence every sentential form has a unique derivation.
\end{proof}

% =====================================================================
\section{The Accountability Type System}
\label{sec:types}
% =====================================================================

\cf{cfc}'s types record \emph{what kind of clinical object} a value is and
\emph{at what floor} it was determined. The system's purpose is to make two
properties statically checkable: that no value claims a resolution finer than
the floor (an accountability property), and that no therapy is declared
\emph{cured} without a localized cell to have closed the gap of.

\subsection{Types and judgements}

\begin{definition}[Types]
\label{def:types}
The \cf{cfc} types are
\[
\tau \bnf \cf{Reading}\alt\cf{Cell}\alt\cf{Coherence}\alt\cf{DiseaseVector}
   \alt\cf{Gap}\alt\cf{Probe}\alt\cf{Panel}\alt\cf{Intervention}
   \alt\cf{Verdict}\alt\cf{Concern}.
\]
\cf{Reading} is a raw measured readout (\Cref{def:probe}); \cf{Cell} is a
localized cell-tight region (\Cref{def:cell}); \cf{Coherence}\,$\in[0,1]$ and
\cf{DiseaseVector}\,$\in[0,1]^{|\Dom|}$ are the health scalars/vectors
(\Cref{def:cell-coherence}); \cf{Gap} is a measured displacement
(\Cref{def:gap}); \cf{Intervention} is a chosen dose vector with its predicted
displacement; \cf{Verdict} is the outcome of a rigidity or coherence admission
(\cf{resolved}/\cf{contested}/\cf{coherent}); \cf{Concern} bundles a cell, its
coherence, and its disease vector (\Cref{rem:concern}). Every value carries a
\emph{floor annotation} $\flo(v)>0$ and, where it denotes a determined region,
a \emph{residue} $\Rd_{\ast}(v)\ge\flo(v)$.
\end{definition}

\begin{definition}[Typing judgement]
\label{def:judgement}
A judgement $\Gamma\vdash e:\tau\,@\,\flo$ reads: ``in floor context $\Gamma$
(fixing the ambient floor and the types of bound identifiers), expression $e$
has type $\tau$ and, if it denotes a determined region, does so at floor
$\flo>0$.'' The context carries the ambient floor set by the most recent
\cf{floor} declaration and the types of bound identifiers.
\end{definition}

\subsection{The floor-positivity rule}

\begin{rulename}[Floor positivity]
\label{rule:positivity}
Every value-introducing rule carries the side condition $\flo>0$ and, for a
region-denoting value, $\Rd_{\ast}\ge\flo$. A literal or expression whose
annotation forces $\flo\le0$ or $\Rd_{\ast}<\flo$ is ill-typed. In particular,
\[
\frac{\Gamma\vdash e:\cf{Panel}\quad\Gamma\vdash r:\cf{Reading}\quad\flo_\Gamma>0}
     {\Gamma\vdash \cf{localize}\ \cdots\ \cf{using}\ e\ \cf{until}\ \cf{resolved}\ \cdots:\cf{Cell}\,@\,\flo_\Gamma}
\]
where $\flo_\Gamma$ is the ambient floor set by the most recent \cf{floor}
declaration.
\end{rulename}

\begin{theorem}[No zero-residue value]
\label{thm:no-zero}
There is no well-typed \cf{cfc} expression of residue $0$: every typeable
region-denoting value (a \cf{Cell}, \cf{Gap}, or \cf{Concern}) has
$\Rd_{\ast}\ge\flo>0$. Equivalently, a point-resolved diagnosis is not
expressible.
\end{theorem}

\begin{proof}
By induction on typing derivations. The only region-introducing rules are the
\cf{localize} rule (\Cref{rule:positivity}, side condition $\flo_\Gamma>0$,
whose result carries residue $=\diam(\Man_S)$ which is $\ge\flo$ by
\Cref{thm:rigidity} at admission), the \cf{close}/\cf{gap} rules (whose gap
residue is the post-contraction distance, $\ge\flo$ by \Cref{thm:convergence}),
and \cf{Concern} construction (which packages a \cf{Cell} and inherits its
residue). Floor-annotated literals carry residue equal to their positive floor
by \Cref{not:floor-literal}. No rule introduces a value of residue $<\flo$, and
$\flo>0$ is a side condition on every floor introduction; no rule lowers the
residue of an already-typed value. Hence no derivation concludes
$\Rd_{\ast}=0$.
\end{proof}

\subsection{The cell-before-cure rule}

\begin{rulename}[Cell before cure]
\label{rule:cell-first}
A \cf{close gap} statement type-checks only if its subject is a \cf{Cell}
(a localized state), never a raw \cf{Reading} or unlocalized region:
\[
\frac{\Gamma\vdash x:\cf{Cell}\,@\,\flo\quad\Gamma\vdash b:\cf{Panel}\ \text{(basis)}}
     {\Gamma\vdash \cf{close}\ \cf{gap}\ \cf{of}\ x\ \cf{toward}\ \cf{coherent}\ \cf{using}\ \cf{basis}\ b\ \cdots:\cf{Intervention}\,@\,\flo}
\]
A \cf{close} whose subject is not of type \cf{Cell} is rejected at compile time.
\end{rulename}

\begin{theorem}[No therapy without localization]
\label{thm:no-blind-therapy}
No well-typed \cf{cfc} program applies a gap-closing intervention to a state
that has not first been localized to a cell. Equivalently, the gap
$\gap(\st)$ of \Cref{def:gap} is well-defined at every typed \cf{close}, because
$\st$ is a \cf{Cell}.
\end{theorem}

\begin{proof}
By \Cref{rule:cell-first} the only rule producing an \cf{Intervention} from a
\cf{close} requires its subject to have type \cf{Cell}. A \cf{Cell} is
introduced only by the \cf{localize} rule (\Cref{rule:positivity}), which runs
the rigidity loop to cell-tightness. Hence any typed \cf{close} acts on a
localized cell, and $\gap$ (defined via the nearest coherent state to a state,
\Cref{def:gap}) is well-defined on it. A \cf{close} on a \cf{Reading} or other
non-\cf{Cell} value has no typing derivation.
\end{proof}

\begin{theorem}[Type soundness]
\label{thm:soundness}
\cf{cfc} typing is sound for the operational semantics of \Cref{sec:opsem}: if
$\Gamma\vdash e:\tau\,@\,\flo$ and $e$ evaluates to a value $v$, then $v$ has
type $\tau$, floor $\flo>0$, and (if $\tau$ is region-denoting) residue
$\Rd_{\ast}(v)\ge\flo$.
\end{theorem}

\begin{proof}[Proof sketch]
Standard syntactic soundness by progress and preservation
\citep{wright1994syntactic,pierce2002types}. \emph{Progress}: a well-typed
closed expression is a value or takes a step, by inspection of the reduction
rules of \Cref{sec:opsem}, each of which has a syntactically disjoint premise
from every other rule of the same constructor. \emph{Preservation}: each
reduction preserves the static type and floor annotation of
\Cref{rule:positivity}, and, by \Cref{thm:no-zero} applied at every step,
preserves $\Rd_{\ast}\ge\flo>0$ for every region value. The only non-standard
element is that reduction is state-mutating (\Cref{sec:opsem}: evaluating
\cf{measure}/\cf{localize}/\cf{close} commits measurements and advances the
committed-measurement count); preservation still holds because the count is
monotone (\Cref{thm:monotone}), so no reduction retroactively lowers a
previously-established residue below its floor. Soundness of the runtime
additionally requires a derived forward model (\Cref{def:soundness}); with a
fitted model the type system remains sound as a syntactic discipline but the
values' correspondence to the patient is only local (\Cref{thm:extrap}).
\end{proof}

% =====================================================================
\section{Operational Semantics}
\label{sec:opsem}
% =====================================================================

We give a small-step semantics \citep{plotkin2004structural} over
configurations $\langle\Sigma,\rec,e\rangle$, where $\Sigma$ is the current
state system with its accumulated committed measurements, $\rec\in\N$ is the
\emph{committed-measurement count} (the program clock), and $e$ is the
expression under evaluation. We write
$\langle\Sigma,\rec,e\rangle\stepto\langle\Sigma',\rec',e'\rangle$ for one
step. Evaluation \emph{is} measurement: a \cf{measure}/\cf{localize}/\cf{close}
step commits measurements into $\Sigma$ and advances $\rec$.

\subsection{Committed measurements and the intrinsic clock}

\begin{definition}[Committed measurement; committed record]
\label{def:committed}
A \emph{committed measurement} is a reduction that administers one probe and
records its reading against $\Sigma$ (refining the surviving set $\Man_S$ or
applying one contraction step). The \emph{committed record} $\rec\in\N$ counts
committed measurements taken so far in the current evaluation.
\end{definition}

\begin{theorem}[Monotonicity of the committed record]
\label{thm:monotone}
Along any evaluation, $\rec$ is non-decreasing and strictly increases at every
committed measurement. No reduction decrements $\rec$; re-evaluating a probe
already administered is a new committed measurement at a strictly higher $\rec$,
never a cached recomputation.
\end{theorem}

\begin{proof}
Every rule administering a probe (\Cref{rule:measure-step,rule:localize-step,%
rule:close-step} below) increments $\rec$ by the number of probes it commits,
which is $\ge1$. No rule of \Cref{sec:opsem} removes a committed measurement or
decrements $\rec$: an ``undo'' (e.g.\ a compensating counter-probe) is itself a
further committed measurement, incrementing $\rec$. Hence $\rec$ is monotone and
strictly increases per committed measurement, and two configurations differing
only in $\rec$ are distinct states.
\end{proof}

\begin{corollary}[Measurement, not simulation]
\label{cor:measure}
Because evaluation mutates $\Sigma$ and advances $\rec$, a \cf{cfc} probe cannot
be a pure function memoised across calls: each occurrence performs the
measurement afresh and is recorded in the clock. Running the program \emph{is}
performing the measurements it names.
\end{corollary}

\subsection{Reduction rules}

\begin{rulename}[Measure commits and counts]
\label{rule:measure-step}
\[
\frac{y=\textsc{Administer}(\Sigma,\Probe_{\mathrm{panel}})}
     {\langle\Sigma,\rec,\cf{measure}\ p\ \cf{yield}\ r\rangle\stepto
      \langle\Sigma\oplus y,\ \rec+\abs{p},\ v_{\cf{Reading}}(y)\rangle}
\]
where $\textsc{Administer}$ integrates the panel's probes against the model and
returns the reading $y$, $\Sigma\oplus y$ records it, and $\abs{p}$ is the
number of probes in the panel.
\end{rulename}

\begin{rulename}[Localize converges]
\label{rule:localize-step}
\[
\frac{\begin{array}{c}
\Man_S=\textsc{Survive}(\Sigma,p)\qquad\diam(\Man_S)<\flo\\
\abs{\text{cells of }\Man_S}=1
\end{array}}
{\langle\Sigma,\rec,\cf{localize}\ x\ \cf{using}\ p\ \cf{until}\ \cf{resolved}\ \cf{yield}\ c\rangle\stepto
 \langle\Sigma',\ \rec+k,\ v_{\cf{Cell}}(\Man_S)\rangle}
\]
where $\textsc{Survive}$ accumulates $k$ probe readings into the surviving set
of \Cref{thm:rigidity} and $v_{\cf{Cell}}$ is the resulting accountable cell.
If instead $\Man_S$ stabilises with more than one cell, the
\cf{otherwise decline} branch fires:
\[
\frac{\textsc{Survive}(\Sigma,p)=\Man_S\text{ stable},\ \abs{\text{cells}}>1}
{\langle\Sigma,\rec,\cf{localize}\ x\ \cf{using}\ p\ \cf{until}\ \cf{resolved}\ \cf{otherwise}\ \cf{decline}\ \cf{yield}\ c\rangle\stepto
 \langle\Sigma',\ \rec+k,\ v_{\cf{decline}}(\Man_S)\rangle}
\]
\end{rulename}

\begin{rulename}[Close contracts to coherence]
\label{rule:close-step}
\[
\frac{\begin{array}{c}
\dose^{\ast}=\textsc{Solve}(\gap(x),\basis_b)\quad x'=x+\textstyle\sum_i \dose^{\ast}_i e_i(x)\\
\dist(x',\Coh)\le\flo
\end{array}}
{\langle\Sigma,\rec,\cf{close}\ \cf{gap}\ \cf{of}\ x\ \cf{toward}\ \cf{coherent}\ \cf{using}\ \cf{basis}\ b\ \cf{until}\ \cf{coherent}\ \cf{yield}\ t\rangle\stepto
 \langle\Sigma',\ \rec+1,\ v_{\cf{Intervention}}(\dose^{\ast})\rangle}
\]
where $\textsc{Solve}$ solves \Cref{prob:closure}. If $\dist(x',\Coh)>\flo$ the
state is re-localized and the rule re-fires on $x'$ (one contraction step per
firing), realising the iteration of \Cref{thm:convergence}; the
\cf{otherwise decline} branch fires if $\textsc{Solve}$ is infeasible
(\Cref{prop:closure-wellposed}).
\end{rulename}

\begin{theorem}[Determinism]
\label{thm:determinism}
For a fixed initial configuration and a fixed forward model and basis, \cf{cfc}
evaluation has at most one normal form: if
$\langle\Sigma,\rec,e\rangle\stepto^{\ast}\langle\cdot,\cdot,v_1\rangle$ and
$\langle\Sigma,\rec,e\rangle\stepto^{\ast}\langle\cdot,\cdot,v_2\rangle$ then
$v_1=v_2$.
\end{theorem}

\begin{proof}
$\textsc{Administer}$, $\textsc{Survive}$, and $\textsc{Solve}$ are total,
deterministic functions of $\Sigma$ and the syntactic clauses (integration of a
fixed model, intersection of enabling sets, and a convex program with a fixed
tie-break, respectively). A \cf{localize} terminates in exactly one of the two
mutually exclusive states of \Cref{rule:localize-step}, determined by whether
$\Man_S$ stabilises at one cell or more---a fact about the fixed model, not a
runtime choice. Two cell-tight surviving sets containing the state are
observationally identical (same cell), so the representative cell is unique. A
\cf{close} likewise deterministically produces the unique minimal-$\ell_1$
displacement under the fixed tie-break, or declines on infeasibility. Hence the
normal form is unique.
\end{proof}

\begin{theorem}[Convergent closure or honest decline]
\label{thm:decline}
Every \cf{localize} over a panel with finitely many registered probes
terminates in exactly one of two states: \textbf{(i)}~\emph{resolved} ---
$\Man_S$ stabilises to a single cell, reported as a \cf{Cell}; or
\textbf{(ii)}~\emph{contested} --- $\Man_S$ stabilises but contains more than
one cell, reported as a \cf{decline} carrying the distinct surviving cells. No
\cf{localize} over a finite panel runs forever without reaching one of these.
\end{theorem}

\begin{proof}
The registered panel is finite, so the reachable set of readings is finite and
$\Man_S$ (an intersection of finitely many enabling constraints) can shrink by
only finitely many steps before every remaining probe has been tried; at that
point $\Man_S$ is stable (rigidity by exhaustion, \Cref{thm:rigidity}). If the
stable $\Man_S$ is cell-tight, case (i) holds; else it meets $>1$ cell and, by
\Cref{cor:cell}, no single cell is licensed, so case (ii) holds and the program
reports the contested cells. The cases are exhaustive and mutually exclusive on
a non-empty surviving set.
\end{proof}

\begin{remark}[Decline is not failure]
\label{rem:decline-not-failure}
\Cref{thm:decline}(ii) formalises the requirement that a clinical runtime be
able to report ``the evidence does not single out one state'' as a first-class,
correctly-terminating outcome---an honest \emph{non liquet}---rather than
looping or silently emitting one contested cell as if unique. A contested
closure is itself useful output: it identifies precisely which cells remain and
what further probe would separate them.
\end{remark}

\begin{theorem}[Scheduler soundness]
\label{thm:scheduler}
The \cf{cfc} scheduler never re-probes a cell already resolved to
cell-tightness, and never starves a convergent \cf{close}: every
gap-contracting step that reduces $\dist(\cdot,\Coh)$ by a fixed factor is
eventually taken, so a contraction (\Cref{def:contraction}) reaches
$\dist\le\flo$ in finitely many steps.
\end{theorem}

\begin{proof}
By \Cref{thm:monotone} the committed record is monotone, so a resolved cell
(recorded in $\Sigma$) is not revisited: re-probing it would strictly increase
$\rec$ without shrinking the already-cell-tight $\Man_S$, and the scheduler,
which advances only on strict progress in $\Man_S$ or $\dist(\cdot,\Coh)$, does
not schedule such a step. For non-starvation: each \cf{close} firing applies one
contraction step reducing $\dist(\cdot,\Coh)$ by a factor $\le\lambda<1$
(\Cref{def:contraction}); by \Cref{thm:convergence} the distance reaches
$\le\flo$ after $k^{\ast}=\lceil\log_{1/\lambda}(\dist_0/\flo)\rceil$ firings,
a finite bound, so the convergent therapy completes.
\end{proof}

% =====================================================================
\section{The Standard Library}
\label{sec:stdlib}
% =====================================================================

Each \cf{cfc} verb binds to a named semantic operation of \Cref{sec:model}--%
\ref{sec:unity}. We state the contract (inputs, output, invoked operation); the
operation's construction is fixed by the semantic model.

\begin{definition}[\textsc{Administer}]
\label{def:administer}
\cf{measure p yield r} invokes $\textsc{Administer}(\Sigma,p)$: integrate each
probe of panel $p$ against the forward model from the current state, returning
the reading $r$ (the species distribution at read time, \Cref{def:probe}). The
reading is a \cf{Reading} value carrying its floor annotation.
\end{definition}

\begin{definition}[\textsc{Survive}, \textsc{Localize}]
\label{def:localize}
\cf{localize x using p until resolved yield c} invokes
$\textsc{Survive}(\Sigma,p)$: accumulate the panel's readings into the surviving
set $\Man_S$ (\Cref{thm:rigidity}) and admit when cell-tight, yielding the
\cf{Cell} $c$; on stable non-cell-tight $\Man_S$ it declines
(\Cref{thm:decline}). The cell carries its coherence $\ecoh$ and disease vector
$\dvec$ (\Cref{cor:coherence-from-cell}).
\end{definition}

\begin{definition}[\textsc{Gap}, \textsc{Solve}, \textsc{Close}]
\label{def:close}
\cf{close gap of x toward coherent using basis b until coherent yield t}
invokes $\textsc{Gap}(x)$ to compute the displacement to the nearest coherent
state (\Cref{def:gap}), then $\textsc{Solve}$ to solve the minimal-$\ell_1$
program (\Cref{prob:closure}) over basis $b$, applies the chosen doses, and
iterates as a contraction (\Cref{thm:convergence}) until the gap is below the
floor, yielding the \cf{Intervention} $t$ (the sparse dose vector with its
predicted displacement). Infeasibility declines
(\Cref{prop:closure-wellposed}).
\end{definition}

\begin{definition}[Soundness of the runtime]
\label{def:soundness}
A \cf{cfc} runtime is \emph{sound} if \textbf{(i)}~$\textsc{Administer}$ agrees
with the true reading up to the uniform derived-model bound of
\Cref{thm:extrap}; \textbf{(ii)}~$\textsc{Survive}$ returns a superset of the
true surviving set (no false exclusion of the actual state);
\textbf{(iii)}~$\textsc{Solve}$ returns a feasible intervention whose simulated
displacement matches its physical displacement up to the same bound. Soundness
reduces to the derived-model guarantee (\Cref{prin:derived}).
\end{definition}

\begin{remark}[The verbs are one primitive]
\label{rem:one-primitive}
$\textsc{Administer}$, $\textsc{Survive}$, and $\textsc{Solve}$ are the single
characterised-perturbation primitive read three ways: as a raw reading, as an
accumulation of readings into a cell, and as a displacement chosen to close a
gap. By \Cref{thm:unity} the measuring and moving faces are the same act; the
standard library is therefore not three independent operations but one act
exposed at three roles. No theorem depends on this reading.
\end{remark}

% =====================================================================
\section{Worked Programs}
\label{sec:examples}
% =====================================================================

\subsection{Localize a patient to a cell}

\begin{lstlisting}[caption={Diagnosis is localization to a cell, never a label.}]
-- localize.cfc
floor 1e-3                 -- ambient floor; nothing resolves finer

basis Metabolic {
  probe glucose_challenge {
    reference: healthy_ref
    observe:   ingest("glucose", 75) >> glycolysis >> tca >> etc
  }
  probe lactate_probe {
    reference: healthy_ref
    observe:   ingest("pyruvate", 10) >> ldh_branch
  }
}

patient P {
  measure panel { glucose_challenge() || lactate_probe() } yield r

  -- run the rigidity loop: keep probing until one cell survives,
  -- otherwise honestly decline (the evidence does not single out a state)
  localize r
    using panel { glucose_challenge() >> tca_followup() }
    until resolved otherwise decline
    yield here

  -- `here : Cell` carries coherence and the disease vector; no disease name
  assert here.coherence >= 0.0#1e-3
  yield here
}
\end{lstlisting}

\subsection{Close the gap without naming a disease}

\begin{lstlisting}[caption={Therapy is gap-closure to the coherent set.}]
-- treat.cfc
floor 1e-3
import Metabolic

patient Q {
  localize baseline
    using panel { glucose_challenge() >> tca_followup() }
    until resolved
    yield here              -- here : Cell (must be a Cell before `close`)

  -- close the measured gap toward coherence with the sparsest set of
  -- characterised interventions, chosen by effect, not by indication
  close gap of here
    toward coherent
    using basis Metabolic
    until coherent otherwise decline
    yield plan             -- plan : Intervention (sparse dose vector)

  emit "gap closed to floor; no disease named"
  yield plan
}
\end{lstlisting}

\subsection{Monitor continuously}

\begin{lstlisting}[caption={A continuously working probe re-reads the state.}]
-- monitor.cfc
floor 1e-3
import Metabolic

patient R {
  monitor R every 3600 {          -- re-run every 3600 time units
    measure panel { lactate_probe() } yield r
    localize r using panel { lactate_probe() } until resolved yield now
    -- raise concern iff coherence falls below threshold (Remark on concern)
    assert now.coherence >= 0.7#1e-3
      emit "cause for concern: coherence below threshold"
  }
}
\end{lstlisting}

\begin{remark}[The same keyword set, three clinical jobs]
\label{rem:three-jobs}
Across the three programs, \cf{measure} reads, \cf{localize} diagnoses (to a
cell, or declines), and \cf{close} treats (to the floor, or declines). No
program names a disease; each returns a cell, an intervention, or an honest
decline. A reader who understands the characterised perturbation understands
the whole language.
\end{remark}

% =====================================================================
\section{Complexity}
\label{sec:complexity}
% =====================================================================

\begin{proposition}[Per-judgment complexity]
\label{prop:complexity}
Let $n=\dim\Man$, $m=\abs{\basis}$, $d=\abs{\Dom}$, and let a cascade have
depth $\ell$ and a reading return $q$ species. Then:
\begin{enumerate}[label=\textup{(\roman*)},leftmargin=2.2em,itemsep=2pt]
\item \cf{measure} costs one numerical integration of an $n$-dimensional system
to the read time, plus $\mathcal{O}(q)$ to read back the species;
\item \cf{localize} costs $\mathcal{O}(q)$ constraint intersections against
precomputed enabling sets per probe, and $\mathcal{O}(d)$ to compute the cell's
coherence and disease vector (\Cref{prop:coherence-bounds});
\item \cf{close} is a projection onto $\Coh$ plus an $\ell_1$-regularised convex
program in $m$ variables and $n$ constraints, solvable in polynomial time
\citep{boyd2004convex};
\item the \cf{close} contraction loop multiplies (iii) by
$\mathcal{O}(\log_{1/\lambda}(1/\flo))$ steps (\Cref{thm:convergence}).
\end{enumerate}
\end{proposition}

\begin{proof}
(i) is one integration and a linear read-back. (ii) each probe adds one
enabling-set intersection over $q$ species; coherence and disease vector are
$\mathcal{O}(d)$ by \Cref{prop:coherence-bounds}. (iii) the gap projection is a
nearest-point computation onto the convex $\Coh$ and the dose solve is a convex
program of the stated size. (iv) \Cref{thm:convergence} bounds the number of
contraction steps by $\lceil\log_{1/\lambda}(\dist_0/\flo)\rceil$, each of cost
(iii).
\end{proof}

% =====================================================================
\section{Constructive Validation}
\label{sec:validation}
% =====================================================================

Because every object is a dynamical system on a finite-dimensional space and
every quantity is an integration, a projection, an intersection, or a sparse
program, each theorem is checkable by direct computation on explicit model
instances, without clinical data. The following protocol validates the calculus
and any conforming implementation \citep{sandve2013ten,wilkinson2016fair}.

\begin{enumerate}[leftmargin=1.6em,itemsep=3pt]
\item \textbf{Floor} (\Cref{thm:floor}). On a model $\Sigma$ with fixed readout
resolution $\delta$, confirm the induced state-discrimination has a positive
infimum $\flo$ and that no state resolves below it.
\item \textbf{Coherence} (\Cref{def:cell-coherence}, \Cref{prop:coherence-bounds}).
Confirm $\ecoh_i,\ecoh\in[0,1]$ at optimal/degraded metric values and that the
coherent set $\Coh_\theta$ is exactly the states with all
$\ecoh_i\ge\theta$.
\item \textbf{Localization} (\Cref{thm:localization}). Sample state pairs;
confirm a floor-injective probe never maps distinct cells to the same reading
within $\delta$, and exhibit a non-injective probe that does, demonstrating
insufficiency.
\item \textbf{Certificate} (\Cref{thm:certificate}). Construct a state-gated
cascade; confirm species $P_k$ appears iff every prior transition is
state-enabled, and that states failing a transition produce $P_k$ absent.
\item \textbf{Rigidity} (\Cref{thm:rigidity}). Grow the observed species set and
confirm $\Man_S$ is non-increasing and becomes cell-tight exactly when no
alternative cell survives.
\item \textbf{Gap-closure} (\Cref{prob:closure}, \Cref{thm:convergence}). On
states off $\Coh$, solve the sparse program and iterate the contraction; confirm
geometric contraction of $\dist(\cdot,\Coh)$ and sparsity of the selected
intervention set.
\item \textbf{Unity} (\Cref{thm:unity}). Confirm a single simulated perturbation
yields both a localizing trace and a state displacement from one integration.
\item \textbf{Decline} (\Cref{thm:decline}). Construct a model with two disjoint
enabling clusters reachable by different probes; confirm \cf{localize} reports
\emph{contested} decline when a probe separating them is withheld and
\emph{resolved} when it is supplied.
\item \textbf{Derived vs.\ fitted} (\Cref{thm:extrap}). Fit a surrogate model on
a training region, then evaluate both surrogate and derived model off-region;
confirm the surrogate's trace error grows off-region while the derived model
stays within its uniform bound---quantifying why free exploration requires the
derived runtime.
\item \textbf{Type safety} (\Cref{thm:no-zero,thm:no-blind-therapy,thm:soundness}).
Confirm the type checker rejects (a) any literal or value of zero residue and
(b) any \cf{close} whose subject is not a \cf{Cell}.
\end{enumerate}

Each check is a finite, deterministic computation whose predicted and measured
outcomes agree exactly for identities and within the stated bound for
inequalities. The suite validates the calculus on explicit model instances; its
fidelity to a real patient is inherited entirely from the fidelity of the
derived forward model, the separate physical question the runtime must answer.

% =====================================================================
\section{Reference Execution Model and Implementation}
\label{sec:impl}
% =====================================================================

A conforming \cf{cfc} implementation must provide, from this specification:

\begin{enumerate}[leftmargin=1.6em,itemsep=2pt]
\item the shared front end: lexer (\Cref{sec:lexical}), parser to the grammar
(\Cref{def:grammar}), and accountability type checker (\Cref{sec:types})
enforcing \Cref{thm:no-zero,thm:no-blind-therapy};
\item the operational semantics of \Cref{sec:opsem}: a state-system
configuration, a monotone committed-measurement count, and the
\cf{measure}/\cf{localize}/\cf{close}/\cf{monitor} reductions, with evaluation
mutating state (\Cref{cor:measure});
\item the standard-library bindings of \Cref{sec:stdlib} to
$\textsc{Administer}$, $\textsc{Survive}$, $\textsc{Solve}$;
\item a \emph{derived} forward model (\Cref{def:fitvsder}) satisfying
runtime soundness (\Cref{def:soundness}); a fitted surrogate may be offered only
for exploration within its training region, and must be flagged as such.
\end{enumerate}

\begin{remark}[Reference interpreter]
\label{rem:reference-impl}
A direct interpreter of \Cref{sec:opsem} over an exact solver---numerical
integration for $\textsc{Administer}$, exact enabling-set intersection for
$\textsc{Survive}$, and a convex-program solver for $\textsc{Solve}$---is the
simplest conforming implementation and serves as the oracle against which
faster or approximate runtimes are checked. A conformance suite must verify the
validation protocol of \Cref{sec:validation} and, in particular, that the type
checker rejects zero-residue values and blind therapies.
\end{remark}

% =====================================================================
\section{Related Paradigms and Scope}
\label{sec:related}
% =====================================================================

\cf{cfc} draws its formal machinery from three standard traditions and departs
from clinical convention in one decisive respect.

From \textbf{programming-language theory} it inherits the discipline of a
lexical grammar, a context-free grammar in EBNF, a static type system with
progress-and-preservation soundness, and a small-step operational semantics
\citep{aho2006compilers,pierce2002types,wright1994syntactic,plotkin2004structural};
its gradual-accountability flavour (every value carrying the floor at which it
was determined) is a domain-specific refinement of the idea that types can
track a quantitative invariant.

From \textbf{systems biology and constraint-based modelling} it inherits the
view of physiology as a kinetic dynamical system and of intervention as
constrained flux redirection \citep{kitano2002systems,palsson2015systems,%
orth2010fba}; the demand that the runtime be derived rather than fitted
(\Cref{sec:runtime}) is the formal counterpart of the modelling community's
preference for mechanistic over regression models when extrapolation matters.

From \textbf{optimization} it inherits the sparse-selection geometry of the
$\ell_1$ program \citep{tibshirani1996lasso,candes2006robust,boyd2004convex} and
the contraction-mapping fixed-point theory \citep{banach1922} that make therapy
a convergent, minimal-intervention procedure.

Its \textbf{departure from convention} is the elimination of the named disease.
Where nosology partitions a continuous pathophysiology into catalogue entries
\citep{feinstein1967clinical,scheuermann2009toward,loscalzo2007network},
\cf{cfc} returns a cell of state space and a measured gap. Naming, if desired,
is a downstream description of which cell was occupied; it is never a
prerequisite for either diagnosis or therapy (\Cref{prin:goal,prin:cure}). The
scope of the language is exactly this: to reduce clinical reasoning to
localization and gap-closure over a derived state, relocating the clinician from
precedent-retrieval to the specification of \emph{intent}---which coherent
target, how far, and at what cost to the particular person---the one input the
calculus does not compute.

% =====================================================================
\section{Conclusion}
\label{sec:conclusion}
% =====================================================================

We have specified Cause-for-Concern as a language whose single primitive is the
characterised perturbation---an administered compound of known composition read
simultaneously for its fate (a measurement) and its displacement (an
intervention)---and whose every operation follows from five proved facts. From a
strictly positive resolution floor we obtained that a state is a cell, so
diagnosis is localization to a cell rather than recognition against a disease
catalogue (\Cref{thm:floor}, \Cref{prin:goal}). Health was defined intrinsically
as membership of a coherent set, quantified by a scalar coherence and a disease
vector, so no nosology is needed to say where health lies
(\Cref{sec:coherence}). A probe localizes the state exactly when its reading map
is floor-injective (\Cref{thm:localization}); a cascading probe strengthens this
to a trajectory certificate, since its later products exist only if the earlier
path occurred, so a false account cannot forge the reading and the correct
stopping rule is rigidity, not a threshold (\Cref{thm:certificate,thm:rigidity}).
Therapy is the closure of the measured gap to the coherent set by the sparsest
combination of characterised interventions selected by effect, executed as a
contraction to a fixed point, with no disease ever named
(\Cref{prob:closure,thm:convergence}). Diagnosis and therapy coincide, because
every characterised perturbation both measures and moves, and the loop halts
precisely when the state is coherent and rigidly known---the operative
definition of wellness (\Cref{thm:unity,thm:wellness}).

On this foundation we built the language in full: a lexical and context-free
grammar; an accountability type system in which no zero-residue value and no
blind therapy are well-typed, proved sound by progress and preservation; a
small-step operational semantics in which evaluation is measurement, with a
monotone committed-measurement clock, proved deterministic and scheduler-sound;
a standard library binding each verb to its semantic operation; a complexity
account; a constructive validation protocol requiring no clinical data; and a
reference execution model. The one substantive demand---that the forward model
be derived, not fitted---was established as a theorem about extrapolation,
because the safety of risk-free \emph{in silico} experimentation reduces to it.
The result is a deterministic, disease-free, report-free calculus of medicine in
which the patient's state is read from a self-certifying physiological program,
the therapeutic target is a measured gap, and the clinician's role is relocated
from precedent-retrieval to judgment over a derived, explorable state. The
specification is complete and implementable as written.

\section*{Data Availability}
No empirical data was used in this specification. The validation protocol of
\Cref{sec:validation} specifies how model instances and, subsequently, clinical
data should be collected and compared to the language's predictions.

\section*{Code Availability}
Implementation requirements are given in \Cref{sec:impl}. A reference
interpreter realises the operational semantics of \Cref{sec:opsem} directly over
an exact solver.

\bibliographystyle{plainnat}
\bibliography{references}

\end{document}
